\documentclass[10pt]{article}
\usepackage[margin=0.8in]{geometry}
\usepackage[T1]{fontenc}
\usepackage{lmodern,amsmath,amssymb,microtype}
\usepackage[colorlinks=true,urlcolor=blue,linkcolor=blue]{hyperref}
\setlength{\parindent}{0pt}
\setlength{\parskip}{4pt}
\title{A once-punctured torus contradicts the finite counting formula\\
\large Disproof of the explicit finite-count clause in problem 00000000568}
\author{lizaixi01}
\date{24 October 2026}
\begin{document}
\maketitle
\vspace{-1em}

\textbf{Claim and scope.}
The problem prints, in both languages, the necessary special case
\[
 N(g,1)=6g-3+\sum_i2^i\binom{2g}{i}.
\]
Under the usual complete binomial sum, this gives $N(1,1)=3+(1+4+4)=12$.
We use the closed oriented torus with one interior marked point (one puncture).
Its coefficient-free cluster algebra has infinitely many distinct cluster variables.
The contradiction concerns this explicit finite formula, not the unspecified general
Poincar\'e-coefficient formula and not a count of variables in a single cluster.

\textbf{The published geometric model and its boundary.}
Fomin--Shapiro--Thurston, Example 4.6, gives the once-punctured torus matrix
\[
 B=\begin{pmatrix}0&2&-2\\-2&0&2\\2&-2&0\end{pmatrix}.
\]
One may obtain the standard ideal triangulation from a fundamental square with
opposite sides identified and a diagonal: there are three arc labels and two
oriented triangular faces, with the same cyclic order of labels on both faces.
Each face contributes one cyclic signed adjacency matrix; their sum is $B$.
The finite adjacency calculation is included in Lean.
Definitions 5.1--5.2 specify seeds, mutations, all cluster variables and their generated
algebra. Theorem 7.11, with its explicit closed-one-puncture exception, uses the
\emph{plain} component here; we do not combine the plain and notched components.

The identification of this exact published matrix model with the marked torus,
and the geometric interpretation of its flips, are the cited external geometric
definition bridge. Lean proves the complete algebraic statement for that exact
model. It does not reprove the classification of surfaces or the arc-complex
comparison theorem, and it introduces no axiom asserting those results.

\textbf{The full seed model.}
Let $P=\mathbb Q[X_0,X_1,X_2]$ and $F=\operatorname{Frac}(P)$, using native
multivariate polynomials and the native fraction field. The initial coordinates
$x_i$ are the actual images of $X_i$ in $F$. They are algebraically independent
and generate $F$ as a field over $\mathbb Q$.
A formal seed consists of its three coordinates, its exchange matrix $B$ or $-B$,
and proved algebraic independence and field generation. The native matrix mutation
formula of Definition 4.7 is implemented, and every mutation sends $B$ to $-B$
and $-B$ to $B$. For any index $k$ (indices modulo $3$), its coordinate mutation is
\[
 u'_k=\frac{u_{k+1}^{\,2}+u_{k+2}^{\,2}}{u_k},\qquad
 u'_i=u_i\quad(i\ne k).
\]
This is proved equal to the usual positive-column monomial plus negative-column
monomial exchange formula for the actual exchange matrix.
All denominators are nonzero by algebraic independence. The numerator is nonzero
because it is the evaluation of a nonzero polynomial; evaluating that polynomial
at $(1,1,1)$ gives $2$.

Every mutation preserves field generation: the unchanged coordinates and the new
coordinate recover the old coordinate by the same exchange equation. A generating
triple in $F$, whose transcendence degree is at least $3$ by the initial independent
triple, is algebraically independent. This is proved using native transcendence
basis results. Mutations are actual involutions, not relations imposed by an axiom.
Consequently every finite word in \emph{all three} mutation directions gives a valid
seed. Let $\mathcal X\subset F$ be precisely the set of coordinates of all these
seeds. The algebra is the actual native $\mathbb Z$-subalgebra
$\mathbb Z[\mathcal X]\subset F$. Thus the coefficient-free convention is used
with integer ground ring; the variable set has not been rationalized or replaced
by word labels. Coefficients in the exchange relation are both $1$, as for the
trivial coefficient semifield.

\textbf{A legal specialization and its invariant.}
Use the native localization
\[
 L=P[(X_0X_1)^{-1}].
\]
It injects into $F$. There is also a genuine ring homomorphism
$e:L\to\mathbb Q$ sending each $X_i$ to $1$, because the sole inverted generator
$X_0X_1$ maps to the unit $1$. No field homomorphism $F\to\mathbb Q$ with these
values is used or claimed.
Write $x,y,z$ for the local coordinate images and set
\[
 c=(x^2+y^2+z^2)(xy)^{-1},\qquad
 f_0=x,\quad f_1=y,\quad f_{n+2}=cf_{n+1}-f_n\quad\text{in }L.
\]
The localization inverse gives $cxy=x^2+y^2+z^2$. Induction proves the genuine
polynomial invariant and exchange identity
\[
 f_n^2+f_{n+1}^2+z^2=cf_nf_{n+1},\qquad
 f_nf_{n+2}=f_{n+1}^2+z^2.
\]
For $a_n=e(f_n)$, these definitions imply
$a_0=a_1=1$ and $a_{n+2}=3a_{n+1}-a_n$.
Induction gives $0<a_{n+1}$ and $a_n\le a_{n+1}$ for all $n$. Hence
$a_{n+2}\ge2a_{n+1}>a_{n+1}$, and $n\mapsto a_{n+1}$ is strictly increasing.
In particular all $f_n$ and their images in $F$ are nonzero.

\textbf{The infinite family consists of actual cluster variables.}
Let $\bar f_n$ denote the image of $f_n$ in $F$. Starting with the initial seed,
repeat mutation in direction $0$ and then direction $1$, a total of $r$ pairs.
The proved exchange identities and nonzero denominators give its exact cluster
\[
 (\bar f_{2r},\ \bar f_{2r+1},\ x_2).
\]
This path is embedded in the full finite-word seed model; direction $2$ is still
available in that model. Thus $\bar f_{2r+1}\in\mathcal X$ for every $r$.
If two of these actual rational functions were equal, injectivity of $L\to F$
would give equality in $L$. Applying $e$ would give
$a_{2r+1}=a_{2s+1}$, so strict increase yields $r=s$.
This is an injection $\mathbb N\hookrightarrow\mathcal X$ and proves native
\texttt{Set.Infinite}. In fact the proof excludes every finite set representing
$\mathcal X$, for every proposed finite cardinal $N$, and therefore excludes $12$.
It never uses the default value of \texttt{Nat.card} on an infinite type.

\textbf{Formal verification and limits.}
The complete proof ends with the following theorem:
\begin{center}\texttt{lean/Main.lean: complete\_counterexample}\end{center}
Use Lean 4.33.0 and the pinned Mathlib revision
\begin{center}\texttt{db584cd6d46c92f209a44c0f1c829460d327499d}.\end{center}
From the \texttt{lean/} directory run the following commands, or run
\texttt{lean --version} followed by \texttt{./verify.sh} for the last three checks:
\begin{verbatim}
lean --version
lake build +Main
lake env lean -DwarningAsError=true Check.lean
lake env leanchecker --fresh Main
\end{verbatim}
All actual namespace declarations, including structure-generated auxiliaries,
are enumerated. Only \texttt{propext}, \texttt{Classical.choice} and
\texttt{Quot.sound} are allowed. The finite adjacency checks use ordinary kernel
decision proofs. There is no \texttt{sorry}, native decision evaluation, or extra
axiom. Exact commands, source hashes, raw logs and execution observations accompany
the submission. Production timings are diagnostic and are not controlled
research-efficiency evidence. The full original bilingual source is retained.

\begin{thebibliography}{1}
\bibitem{FST}
S. Fomin, M. Shapiro and D. Thurston,
\emph{Cluster algebras and triangulated surfaces. Part I: Cluster complexes},
Acta Mathematica \textbf{201} (2008), 83--146.
\href{https://arxiv.org/pdf/math/0608367v3}{arXiv:math/0608367v3},
printed pp.17--18 (Example 4.6, Definitions 4.7 and 5.1--5.2) and p.28
(Proposition 7.10 and Theorem 7.11). These selected pages were read; no whole-paper
read is claimed.
\end{thebibliography}
\end{document}
